% META: {"id": "TSPE-GEOESPACE-CR-015", "chapitre": "TSPE-GEOMETRIE-ESPACE", "type_objet": "cours", "capacites_codes": ["C8", "C11", "C14"], "sources_inspiration": [], "mode_creation": "ex_nihilo", "status": "approved"}

\section{Projection orthogonale et distances}

% ============================================================
% STRATE 1 — Demonstration exigible et methode de calcul
% ============================================================

\definition[1]{
Le \textbf{projete orthogonal} d'un point $M$ sur un plan $\mathcal{P}$ est le point $H$, intersection de $\mathcal{P}$ avec la droite passant par $M$ et de vecteur directeur un vecteur normal a $\mathcal{P}$.
}

\theoreme{
Le projete orthogonal $H$ de $M$ sur $\mathcal{P}$ est le point de $\mathcal{P}$ le plus proche de $M$ : pour tout point $N \in \mathcal{P}$, $N 
eq H$, on a $MH < MN$.
}

\demonstration{
Soit $N \in \mathcal{P}$, $N 
eq H$. Comme $H,N \in \mathcal{P}$, le vecteur $\vec{HN}$ est un vecteur du plan $\mathcal{P}$.

Par construction, $(MH)$ est orthogonale a $\mathcal{P}$, donc $\vec{MH}$ est orthogonal a tout vecteur de $\mathcal{P}$, en particulier a $\vec{HN}$ : $\vec{MH}\cdot\vec{HN}=0$.

Par la relation de Chasles, $\vec{MN} = \vec{MH}+\vec{HN}$, donc dans le triangle $MHN$ (rectangle en $H$ d'apres ce qui precede), le theoreme de Pythagore donne :
\[
  MN^2 = MH^2+HN^2.
\]

Comme $N 
eq H$, $HN>0$, donc $HN^2>0$, d'ou $MN^2 = MH^2+HN^2 > MH^2$, soit $MN>MH$ (les deux membres sont positifs).

Ainsi $H$ est strictement plus proche de $M$ que tout autre point de $\mathcal{P}$.
}

\margeAppui{
Cette demonstration repose sur le theoreme de Pythagore applique au triangle rectangle $MHN$, consequence directe de l'orthogonalite de $(MH)$ et de $\mathcal{P}$.
}

\propriete[Distance d'un point à un plan]{
$d(M,\mathcal{P}) = MH$, ou $H$ est le projete orthogonal de $M$ sur $\mathcal{P}$. Si $\mathcal{P}$ a pour equation $ax+by+cz+d=0$ (repere orthonorme) et $M(x_0,y_0,z_0)$ :
\[
  d(M,\mathcal{P}) = \frac{|ax_0+by_0+cz_0+d|}{\sqrt{a^2+b^2+c^2}}.
\]
}

\exemple{
Calculer les coordonnees du projete orthogonal de $M(1,2,3)$ sur le plan $\mathcal{P}: x+y+z-3=0$.

Vecteur normal $\vec{n}\begin{pmatrix}1\\1\\1\end{pmatrix}$. La droite $(MH)$ a pour representation parametrique $\begin{cases}x=1+t\\y=2+t\\z=3+t\end{cases}$.

$H \in \mathcal{P}$ : $(1+t)+(2+t)+(3+t)-3=0 \iff 3t+3=0 \iff t=-1$.

Donc $H(1-1,\,2-1,\,3-1) = H(0,1,2)$.

% BEGIN-VERIFY
% t = -1
% H = (1+t, 2+t, 3+t)
% assert H == (0, 1, 2)
% assert H[0]+H[1]+H[2]-3 == 0
% END-VERIFY
}

% ============================================================
% STRATE 2 — Erreurs frequentes
% ============================================================

\erreurFrequente{
\textbf{Oublier de verifier que $H$ appartient bien au plan.} Apres avoir determine $t$ et calcule les coordonnees de $H$, il est recommande de verifier que $H$ verifie l'equation du plan (controle de coherence).
}
